\documentclass[a4paper,12pt]{article}
\usepackage[T2A]{fontenc}
\usepackage[utf8]{inputenc}
\usepackage[russian]{babel}
\usepackage{helvet}
\renewcommand{\familydefault}{\sfdefault}
\usepackage{amsmath,amssymb,mathtools,icomma}
\usepackage[margin=19mm,headheight=15pt,headsep=8mm,footskip=11mm]{geometry}
\usepackage{microtype,xcolor}
\usepackage{tikz}
\usetikzlibrary{calc,arrows.meta,positioning,shapes.geometric}
\usepackage{tkz-euclide}
\usepackage{pgfplots}
\pgfplotsset{compat=1.18}
\usepackage{booktabs,tabularx,enumitem,fancyhdr,setspace,needspace,hyperref}
\definecolor{accent}{HTML}{008080}
\definecolor{darkaccent}{HTML}{004D4D}
\definecolor{textgray}{HTML}{333333}
\definecolor{pale}{HTML}{E8F1EF}
\hypersetup{colorlinks=false,pdfauthor={Лёвин Артём Александрович},pdftitle={Математика ЕГЭ. Чек-лист, 24.10.2026}}
\onehalfspacing
\setlength{\parindent}{0pt}
\setlength{\parskip}{3pt}
\setlist[itemize]{itemsep=2pt,topsep=4pt,leftmargin=1.5em}
\setlist[enumerate]{itemsep=8pt,topsep=7pt,leftmargin=1.65em}
\pagestyle{fancy}
\fancyhf{}
\fancyhead[L]{\small\color{textgray} ЕГЭ / Математика}
\fancyhead[R]{\small\color{textgray} 24.10.2026}
\fancyfoot[C]{\small\color{textgray}\thepage}
\renewcommand{\headrulewidth}{0.25pt}
\newcommand{\ck}{\(\square\)\enspace}
\newcommand{\topic}[1]{\needspace{5\baselineskip}\subsection*{\color{darkaccent}#1}}
\newcommand{\tip}[1]{\par\textcolor{darkaccent}{\textbf{Проверьте:}} #1\par}
\begin{document}
\thispagestyle{empty}
\vspace*{31mm}
\begin{center}
{\Large\color{darkaccent} СИСТЕМА ПОВТОРЕНИЯ ЕГЭ}\par
\vspace{22mm}
{\Huge\bfseries ЧЕК-ЛИСТ}\par
\vspace{8mm}
{\LARGE Математика. Повторение}\par
\vspace{5mm}
{\large Геометрия, векторы, вероятность,\par степени, логарифмы и текстовые задачи}\par
\vspace{27mm}
{\color{accent}\rule{0.63\linewidth}{0.8pt}}\par
\vspace{12mm}
{\Large 24 октября 2026 года}\par
\vspace{10mm}
{\large Путь А \quad $\cdot$ \quad 10 заданий}\par
\vspace{25mm}
{\normalsize Лёвин Артём Александрович}\par
{\bfseries эксклюзивно для Тимофея Васильченко}\par
\end{center}
\vfill
\begin{center}
\begin{tikzpicture}
\draw[accent!65,line width=1.2pt] (0,0) .. controls (2.6,0.7) and (4.5,-0.65) .. (7.2,0) .. controls (9.1,0.44) and (10.4,0.13) .. (11.3,-0.13);
\end{tikzpicture}
\end{center}
\clearpage

\section*{\color{darkaccent}Математика: опорные методы}
\topic{1. Площадь треугольника}
\ck Выбираю формулу по известным величинам:
\[
S=\frac12 ah_a,\qquad S=\frac12 ab\sin\gamma,\qquad
S=\sqrt{p(p-a)(p-b)(p-c)},\quad p=\frac{a+b+c}{2}.
\]
\ck Помню связи с окружностями: $S=pr$ и $S=\dfrac{abc}{4R}$, где $r$ и $R$ --- радиусы вписанной и описанной окружностей.
\par\medskip
\begin{minipage}[c]{0.54\linewidth}
\textbf{Пример.} Боковые стороны равны $\sqrt6$, угол между ними $30^\circ$.
\[
\begin{aligned}
S&=\frac12\cdot\sqrt6\cdot\sqrt6\cdot\sin30^\circ\\
 &=\frac12\cdot6\cdot\frac12=\boxed{\frac32}.
\end{aligned}
\]
\end{minipage}\hfill
\begin{minipage}[c]{0.40\linewidth}
\begin{center}
\begin{tikzpicture}[scale=1.06]
\tkzDefPoint(-0.65,0){A}
\tkzDefPoint(0.65,0){B}
\tkzDefPoint(0,2.426){C}
\tkzDrawPolygon[thick,color=darkaccent](A,B,C)
\tkzMarkSegments[mark=|,color=accent](A,C B,C)
\tkzMarkAngle[size=0.36,color=accent](A,C,B)
\tkzLabelAngle[pos=0.59](A,C,B){$30^\circ$}
\tkzLabelPoints[below](A,B)
\tkzLabelPoints[above](C)
\tkzLabelSegment[left](A,C){$\sqrt6$}
\tkzLabelSegment[right](B,C){$\sqrt6$}
\end{tikzpicture}
\end{center}
\end{minipage}
\par
\tip{В формуле $\tfrac12 ab\sin\gamma$ угол $\gamma$ лежит между сторонами $a$ и $b$.}

\topic{2. Векторы на координатной плоскости}
\ck Умножаю вектор на число покоординатно: $k(x;y)=(kx;ky)$.

\ck Скалярное произведение вычисляю по формуле:
\[
(a_x;a_y)\cdot(b_x;b_y)=a_xb_x+a_yb_y.
\]
\ck Условие перпендикулярности ненулевых векторов: $\vec a\cdot\vec b=0$.

\textbf{Пример.} Если $\vec u=(t;5)$, $\vec v=(4;-8)$, то
\[
\vec u\perp\vec v\iff 4t+5(-8)=0\iff \boxed{t=10}.
\]
\tip{Не теряйте минус в координате; результат скалярного произведения --- число.}

\topic{3. Объёмы шара и цилиндра}
\ck Использую $V_{\text{ш}}=\dfrac43\pi R^3$ и $V_{\text{ц}}=\pi R^2H$.

\ck Если шар вписан в цилиндр, то $H=2R$. Следовательно,
\[
\frac{V_{\text{ш}}}{V_{\text{ц}}}
=\frac{\tfrac43\pi R^3}{2\pi R^3}
=\boxed{\frac23}.
\]
\textbf{Пример.} При объёме цилиндра $180$ объём вписанного шара равен $120$.

\topic{4. Вероятности и таблица возможных исходов}
\ck Для независимых событий $A$ и $B$:
$P(A\cap B)=P(A)P(B)$.

\ck Если случаи несовместны, складываю их вероятности.

\textbf{Пример.} Команда проводит два независимых матча. Дома: победа с вероятностью $0{,}6$, ничья --- $0{,}1$; в гостях: победа --- $0{,}3$, ничья --- $0{,}2$. Победа даёт 3 очка, ничья --- 1. Вероятность набрать не менее 4 очков:
\[
\underbrace{0{,}6\cdot0{,}3}_{3+3}
+\underbrace{0{,}6\cdot0{,}2}_{3+1}
+\underbrace{0{,}1\cdot0{,}3}_{1+3}
=\boxed{0{,}33}.
\]
\tip{Перечислите все подходящие пары результатов: так проще избежать пропуска случая.}

\topic{5. Степени и показательные уравнения}
\ck Привожу обе стороны к одному основанию:
\[
(a^m)^n=a^{mn},\qquad a^{u}=a^{v}\iff u=v
\quad (a>0,\ a\ne1).
\]
\ck Корни представляю как дробные степени: $\sqrt[n]{a}=a^{1/n}$ при $a>0$.

\textbf{Пример.}
\[
4^{3-x}=\sqrt8\iff 2^{6-2x}=2^{3/2}
\iff 6-2x=\frac32\iff \boxed{x=\frac94}.
\]
\tip{При возведении степени в степень перемножаются показатели, а не основания.}
\clearpage

\enlargethispage{2\baselineskip}
\topic{6. Логарифмы: свойства и ограничения}
\ck Проверяю условия: $a>0$, $a\ne1$, аргумент логарифма положителен.
\[
\log_a(uv)=\log_a u+\log_a v,\quad
\log_a u^k=k\log_a u,\quad
\log_a b=\frac1{\log_b a}
\]
для положительных аргументов и допустимых оснований.

\textbf{Пример.} При $\log_b a=-2$:
\[
\log_a(b^6\sqrt a)
=6\log_a b+\frac12
=\frac6{-2}+\frac12=\boxed{-\frac52}.
\]
\ck Если логарифмическое уравнение сводится к алгебраическому, проверяю полученные корни по первоначальной области определения.

\textbf{Пример.}
\[
\log_2(x-1)+\log_2(x+1)=3,\quad x>1,
\]
\[
\log_2(x^2-1)=3\iff x^2-1=8\iff x=\pm3.
\]
Из условия $x>1$ подходит только $\boxed{x=3}$.

\topic{7. Логарифмическая функция и коэффициент}
\ck Для $y=k\log_a x$ помню: $x>0$ и график проходит через точку $(1;0)$.

\ck Если известна точка $(x_0;y_0)$, то $k=\dfrac{y_0}{\log_a x_0}$ при $x_0>0$, $x_0\ne1$.

\begin{minipage}[c]{0.53\linewidth}
\textbf{Пример.} График $y=k\log_2 x$ проходит через $(4;6)$:
\[
6=k\log_2 4=2k\quad\Longrightarrow\quad\boxed{k=3}.
\]
При $a>1$ знак $k$ определяет направление монотонности: $k>0$ --- возрастание, $k<0$ --- убывание.
\end{minipage}\hfill
\begin{minipage}[c]{0.43\linewidth}
\begin{tikzpicture}
\begin{axis}[width=\linewidth,height=39mm,xmin=0,xmax=5,ymin=-5,ymax=8,
 axis lines=middle,xtick={1,2,4},ytick={0,3,6},
 xlabel={$x$},ylabel={$y$},tick label style={font=\scriptsize},
 label style={font=\small},clip=false,samples=200]
\addplot[accent,very thick,domain=0.2:5] {3*ln(x)/ln(2)};
\addplot[only marks,mark=*,darkaccent,mark size=1.8pt] coordinates {(1,0) (4,6)};
\node[anchor=south east,font=\scriptsize] at (axis cs:4,6) {$(4;6)$};
\end{axis}
\end{tikzpicture}
\end{minipage}

\topic{8. Текстовые задачи: выбор переменной и уравнения}
\ck Задаю искомую величину $x$, записываю ограничения на неё и выражаю остальные величины через $x$.

\ck В задачах на путь и время использую тождество $t=S/v$ (при $v>0$). Оно помогает составить уравнение по разнице времён.

\textbf{Пример.} Два участника проехали по $240$ км со скоростями $x$ и $x-12$ км/ч. Второй прибыл на $40$ минут позже:
\[
\frac{240}{x-12}-\frac{240}{x}=\frac23,\quad x>12.
\]
Отсюда $x(x-12)=4320$, то есть $x\in\{72,-60\}$. По условию $x>12$, поэтому $\boxed{x=72\text{ км/ч}}$.

\clearpage
\section*{\color{darkaccent}Алгоритм: вероятность составного события}
\vspace{8mm}
\centering
\tikzset{
start/.style={ellipse,draw=darkaccent,very thick,fill=pale,minimum width=33mm,minimum height=12mm,align=center},
act/.style={rounded corners=2mm,rectangle,draw=accent,very thick,align=center,text width=105mm,minimum height=14mm,inner sep=3mm},
condition/.style={diamond,aspect=2.4,draw=darkaccent,very thick,align=center,text width=44mm,inner sep=1.5mm},
branch/.style={rounded corners=2mm,rectangle,draw=accent,very thick,align=center,text width=42mm,minimum height=16mm,inner sep=2mm},
flow/.style={-{Stealth[length=3mm]},thick,color=darkaccent}}
\begin{tikzpicture}[node distance=10mm and 7mm]
\node[start] (beg) {Начало};
\node[act,below=of beg] (mark) {Перечислите исходы первого и второго испытания};
\node[act,below=of mark] (pairs) {Составьте все возможные пары исходов};
\node[condition,below=13mm of pairs] (good) {Пара удовлетворяет условию задачи?};
\node[branch,below=18mm of good,xshift=-34mm] (yes) {Да: найдите\ вероятность пары};
\node[branch,below=18mm of good,xshift=34mm] (no) {Нет: исключите\ эту пару};
\node[act,below=13mm of yes,xshift=34mm] (sum) {Сложите вероятности всех подходящих пар};
\node[act,below=of sum] (check) {Убедитесь, что $0\le P\le1$ и не пропущены случаи};
\node[start,below=of check] (end) {Ответ};
\draw[flow] (beg)--(mark); \draw[flow] (mark)--(pairs);
\draw[flow] (pairs)--(good);
\draw[flow] (good.south west)--node[sloped,above,font=\small]{да}(yes.north);
\draw[flow] (good.south east)--node[sloped,above,font=\small]{нет}(no.north);
\draw[flow] (yes.south)--([xshift=-25mm]sum.north);
\draw[flow] (no.south)--([xshift=25mm]sum.north);
\draw[flow] (sum)--(check); \draw[flow] (check)--(end);
\end{tikzpicture}
\clearpage

\raggedright
\section*{\color{darkaccent}Тренировочный блок --- путь А}
\textbf{10 задач по математике.} Решайте по порядку: от применения формул к задачам, где нужно самостоятельно выбрать метод.
\begin{enumerate}[label=\textbf{\arabic*.}]
\item В равнобедренном треугольнике боковая сторона равна $2\sqrt3$, а угол при вершине равен $30^\circ$. Найдите площадь.
\item Даны векторы $\vec a=(2;-3)$ и $\vec b=(-4;6)$. Найдите скалярное произведение векторов $2\vec a$ и $0{,}5\vec b$.
\item Шар вписан в цилиндр объёмом $270$ см$^3$. Найдите объём шара в см$^3$.
\item Решите уравнение $9^{2-x}=\sqrt{27}$.
\item Найдите $\log_a(b^4\sqrt a)$, если $\log_b a=-4$ (все логарифмы определены).
\item График функции $y=k\log_2 x$ проходит через точку $(4;-6)$. Найдите $k$.
\item Команда проводит два независимых матча. Дома вероятность победы $0{,}5$, ничьей $0{,}2$; в гостях вероятность победы $0{,}4$, ничьей $0{,}3$. За победу дают 3 очка, за ничью --- 1 очко. Найдите вероятность набрать не менее 4 очков.
\item Решите уравнение $\log_2(x-1)+\log_2(x+1)=3$.
\item Мотоциклист и автомобилист преодолели по $150$ км. Скорость мотоциклиста на $10$ км/ч больше, поэтому он прибыл на $30$ минут раньше. Найдите скорость мотоциклиста.
\item Стороны треугольника равны $13$, $14$ и $15$. Найдите радиус вписанной в него окружности.
\end{enumerate}
\clearpage

\section*{\color{darkaccent}Ответы для самопроверки}
\renewcommand{\arraystretch}{1.55}
\begin{center}
\begin{tabularx}{0.96\linewidth}{@{}>{\bfseries}c X@{}}
\toprule
Номер & Ответ \\
\midrule
1 & $3$ \\
2 & $-26$ \\
3 & $180$ см$^3$ \\
4 & $x=\dfrac54$ \\
5 & $-\dfrac12$ \\
6 & $-3$ \\
7 & $0{,}43$ \\
8 & $x=3$ \\
9 & $60$ км/ч \\
10 & $4$ \\
\bottomrule
\end{tabularx}
\end{center}
\vspace{10mm}
\textbf{На что обратить внимание при проверке}\par
\begin{itemize}
\item Для площади через две стороны нужен угол \emph{между} этими сторонами.
\item В скалярном произведении внимательно отслеживайте отрицательные координаты.
\item В показательных уравнениях корректно приводите степени к общему основанию.
\item В логарифмах проверяйте область определения ещё до преобразований.
\item В задачах на вероятность учитывайте все подходящие несовместные случаи.
\item В текстовых задачах обязательно проверяйте ограничение на найденную переменную.
\end{itemize}
\vfill
{\small\color{textgray}Лёвин Артём Александрович эксклюзивно для Тимофея Васильченко.}
\end{document}
